\section{Conclusion}
\label{sec:conclusion}

We began by asking whether pretrained weight geometry could predict optimal per-layer LoRA rank without any training data. For BERT-base-uncased, it does --- with strong predictive discrimination ($r = 0.984$, $n=5$, bimodal oracle structure).

The effective rank of a pretrained weight matrix, computed in a single SVD pass from $W_0$ alone, correlates with PARA oracle ranks at Pearson $r = 0.984$ ($p = 0.0013$, $n=5$ oracle layers). FFN intermediate matrices (erank $\sim 704$--$710$) receive oracle rank $r^* = 64$; attention layers --- including both output projections and query projections (erank $\sim 557$--$590$) --- receive oracle rank $r^* = 4$. The erank metric cleanly separates these structural tiers without observing a single training example. Participation ratio agrees with erank at $\rho = 0.968$.

Our main contributions are: (1) the first empirical test of a purely structural per-layer LoRA rank predictor, demonstrating strong correlation with PARA oracle ranks for BERT-base-uncased; (2) an open PARA oracle + erank correlation protocol for replication; and (3) erank maps for three transformer families showing consistent layer-type structure.

Three future directions trace to the experimental evidence: (1) partial correlation controlling for depth to separate erank signal from depth confound (see Limitation L6); (2) full oracle sweeps for DeBERTa-v3-base and ViT-base to establish multi-family evidence; and (3) extension to decoder-only LLMs, where ViT's cross-modal result suggests the mechanism may generalize.

The central finding --- that a pretrained weight's singular value entropy predicts its adaptation rank need --- connects pre-training geometry to fine-tuning efficiency in a way that no training signal could achieve. We hope this work encourages broader investigation of what fine-tuning hyperparameters can be predicted from the geometry of pretrained weights alone.
